\subsection{Tourism Components}

Goeldner and Ritchie~\cite{goeldner2009} organise the tourism phenomenon into four interrelated components: the tourist, natural resources and environment, the built environment, and the operating sectors of the tourism industry (Figure~1.2, p.~14). Understanding where hotels sit within this structure clarifies the context of our business problem.

\paragraph{The Tourist.}
``The very heart of the tourism phenomenon model is unequivocally the tourist and the travel experiences that he or she seeks'' (p.~15). Every downstream facility, service, and pricing decision exists to serve the tourist's needs. In the context of revenue management, each tourist is a decision maker who chooses whether to book a room, at what price, at what time, and through which channel. According to AH\&LA figures cited in the textbook, 44\% of lodging customers are business travellers (average nightly rate \$112) and 56\% are leisure travellers (\$103 per night) (p.~156). This difference in willingness-to-pay between segments is the economic foundation for differentiated pricing.

\paragraph{Natural Resources and Environment.}
The physiography, climate, and culture of a destination are immutable---they cannot be changed to increase visitor arrivals. This immutability is the root cause of \emph{seasonality}, one of the five variables in the demand function $D = f(\text{propensity}, \text{resistance})$ that the authors develop in Chapter~13 (p.~414). Seasonality means that demand for hotel rooms fluctuates predictably across the year while room supply remains fixed, forcing managers to rely on price as the primary lever.

\paragraph{The Built Environment.}
The built environment distinguishes \emph{infrastructure} (roads, sewage systems, commercial networks built for residents) from the \emph{tourism superstructure}---``facilities that have been developed especially to respond to the demands of visitors,'' with hotels listed as the first example (p.~17). Crucially, ``the characteristics of components of the superstructure are essentially determined by visitor wishes rather than resident desires'' (p.~17). A hotel's pricing, room types, and service offerings must therefore reflect what visitors want and are willing to pay for.

\paragraph{Operating Sectors.}
The operating sectors---transportation, accommodation, food services, attractions, events, and others---represent what the general public perceives as ``tourism'' (p.~19). Within these sectors, accommodation is the one most directly connected to our business problem, and it is the focus of Chapter~6 of the textbook.

\paragraph{Technology and Information.}
Goeldner and Ritchie single out two recent dimensions of the built environment that have reshaped tourism. On technology: ``The advent of jet aircraft and the massive invasion of telecommunications technology, linked closely with computer technology, has had a dramatic impact on the very essence of the tourism phenomenon'' (p.~18). On information: ``Increasingly, the success of a destination is determined by its ability to assemble, interpret, and utilize information in an effective manner'' (p.~18). This latter statement is the intellectual foundation for our decision support system: the ability to assemble and interpret booking data is what separates a data-driven hotel from one that prices rooms by intuition.
